\section{Decision-Level Epistemic Integrity}
\label{sec:definition}
Prior work formalizes what an agent may \emph{do}: which tools it may
invoke, which flows may reach which sinks~\cite{camel,cxi,argus}. We
formalize what an agent may \emph{conclude}: the conditions under which
the beliefs that support a consequential action can be trusted, and
which of those conditions a runtime monitor can actually enforce. Every
object the invariants quantify over is declared once, in
Definition~\ref{def:policy}, before any invariant uses it; the running
example of \S\ref{sec:setting} is instantiated as soon as the objects
exist and is revisited after each result.

\subsection{Policy Objects and Decision Traces}
\label{sec:prelim}

\begin{definition}[Integrity lattice, labels, and policy objects]
\label{def:policy}
Let $(L, \sqsubseteq, \sqcap, \sqcup)$ be a finite lattice of integrity
labels with top $\top$ (maximally trusted) and bottom $\bot$
(untrusted), in the sense of Biba~\cite{biba} and
Denning~\cite{denning}. Let $\mathit{Ch}$ be the finite set of ingress
channels (user turn, tool result, retrieved document, memory read, model
output), $C$ the set of action classes, and $\mathit{Slots}$ the set of
named argument positions. A \emph{policy} $\Pi$ fixes:
\begin{itemize}[leftmargin=1.15em,itemsep=1pt]
\item an \emph{ingress labeling} $\lambda : \mathit{Ch} \to L$ and a
\emph{floor} $\lmin \in L$; model-generated and external content default
to low labels;
\item \emph{protected classes} $\Cprot \subseteq C$ (irreversible,
external, or safety-relevant effects), a \emph{required label}
$\lreq : \Cprot \to L$, and a \emph{sufficiency threshold}
$\tau : \Cprot \to [0,1]$ (written $\tc$ for $\tau(c)$);
\item a \emph{binding policy} $\bindpol$ assigning to each
$c \in \Cprot$ a set $\bindpol(c) \subseteq \mathit{Slots}$ of bound
argument slots and, for each $s \in \bindpol(c)$, a required label
$\ell_{\bindpol}(c, s) \in L$;
\item an \emph{endorsement policy} $\pi$, a finite set of authorized
label-raising rules (e.g., $k$-source corroboration, verifier-tool
confirmation, human sign-off), each with a predicate
$\mathsf{Authorized}(\rho, b, \ell')$ stating when rule $\rho \in \pi$
may raise belief $b$ to label $\ell'$;
\item a \emph{tolerance} $\varepsilon : \Cprot \to \mathbb{R}_{\ge 0}$
and a class $\Dirrel(q)$ of task-irrelevant perturbations
(Def.~\ref{def:i5}).
\end{itemize}
Every evidence atom and belief carries a label $\ell(\cdot) \in L$;
for an atom $e$ ingested on channel $\mathit{ch}$, $\ell(e) :=
\lambda(\mathit{ch})$. The concrete lattice used throughout the paper
is $\mathrm{UNTRUSTED} \sqsubset \mathrm{MEDIUM} \sqsubset
\mathrm{TRUSTED}$, with $\lambda(\text{user turn}) = \mathrm{TRUSTED}$
(session-signed), $\lambda(\text{tool result}) = \lambda(\text{memory})
= \mathrm{MEDIUM}$ (channel-attested), and $\lambda(\text{model
output}) = \mathrm{UNTRUSTED}$; the deployed values of the remaining
objects are given in \S\ref{sec:evalsetup}.
\end{definition}

\begin{definition}[Evidence and attestation]
\label{def:evidence}
An \emph{evidence atom} $e \in E$ is an immutable record of ingested
content with label $\ell(e)$, channel $\mathit{ch}(e)$, source metadata,
and an optional attestation. The predicate $\att(e)$ holds iff $e$
carries a signature chain that verifies against a trusted root and
whose validity window covers the current time. The freshness clause
subsumes long-horizon staleness: expired attestations fail $\att$, so
drift over long deployments surfaces as a provenance-completeness
violation rather than a silent decay.
\end{definition}

\paragraph*{What an attestation is, concretely}
In CZTA (\S\ref{sec:czta}) an attestation is the tuple (channel
credential, content digest, issue time, validity window, signature).
For user turns the credential is the session's RS256 token, verified
against the identity provider's published JWKS; for tool results and
memory reads it is the gateway's Ed25519 receipt over the content
digest, issued at ingress and chained into the ledger. The trusted root
is thus the realm key for user channels and the ledger key otherwise;
windows are the token expiry and a per-channel time-to-live; $\att$
fails on a bad signature, an unknown signer, a digest mismatch, or a
lapsed window. The definition does not depend on this choice; it is
stated so that the TCB of \S\ref{sec:tcbdef} is concrete.

\begin{definition}[Beliefs, provenance, updates, endorsements]
\label{def:beliefs}
A \emph{belief} $b \in B$ is a labeled assertion the agent system holds
at runtime, with dependency set $\deps(b) \subseteq E \cup B$ and a
finite set $\vals(b) \subseteq V$ of the values it names (identifiers,
principals, amounts, file names), extracted at creation time by the
recorder from the belief's content. The \emph{provenance graph} $P$ is
the directed acyclic graph whose nodes are $E \cup B$ and whose edges
run from each $b$ to every element of $\deps(b)$; it is recorded in an
append-only, hash-chained log (we use the W3C~PROV vocabulary for the
edge kinds~\cite{prov}, a recording convention that carries no semantic
commitment). Two event kinds modify beliefs:
\begin{itemize}[leftmargin=1.15em,itemsep=1pt]
\item an \emph{update} $\ev{upd}(b', X)$ creates $b'$ with
$\deps(b') = X$ and default label
$\ell(b') := \bigsqcap_{x \in X} \ell(x)$. This is the Biba derivation
rule~\cite{biba}; its guarantee is that a derived belief's label is a
certified lower bound on the trust of everything it was derived from,
so no derivation can manufacture trust;
\item an \emph{endorsement} $\ev{end}(b, \ell', \rho)$ raises
$\ell(b)$ to $\ell'$ under a rule $\rho \in \pi$ with
$\mathsf{Authorized}(\rho, b, \ell')$, in the sense of endorsement in
decentralized IFC~\cite{myersliskov}. Endorsements are logged,
attributable, and rate-limitable; they are the only lawful label
increases.
\end{itemize}
\end{definition}

\begin{definition}[Actions, arguments, and decision traces]
\label{def:actions}
An \emph{action event} $\ev{act}(a)$ has class $c(a) \in C$, an
\emph{argument map} $\args(a) : \mathit{Slots} \rightharpoonup V$, and a
\emph{support set} $\support(a) \subseteq B$, the beliefs the system
commits to as justifying $a$. Fix the event alphabet
$\Sigma = \{ \ev{ing}(e, ch),\, \ev{upd}(b, X),\, \ev{end}(b, \ell, \rho),\, \ev{act}(a) \}$.
A \emph{decision trace} is a finite sequence
$\sigma = \epsilon_1 \epsilon_2 \cdots \epsilon_k$ over $\Sigma$;
$\acts(\sigma)$ denotes its protected action events. For a task input
$q$ and untrusted-channel content $d$, we write $J(q, d)$ for the
protected action the system selects at a designated decision point,
holding trusted state fixed.
\end{definition}

\paragraph*{Epistemic TCB}
\label{sec:tcbdef}
The definition is stated relative to a trusted computing base: the
ingress labeler $\lambda$ together with the channel authentication it
relies on, the attestation verifier and its keys, the append-only
provenance recorder including the extraction of $\vals(\cdot)$, the
support aggregator $S$ (Def.~\ref{def:i3}), and the policy engine and
gateway. The model, all tool outputs, retrieved content, and memory
contents are \emph{untrusted}. Making the TCB explicit answers the
observation that provenance assignment is the under-specified trusted
base of out-of-band defenses~\cite{oob}; we do not eliminate that base,
we name it, minimize it, and measure it (\S\ref{sec:tcb}). The labeler
is deliberately simple: it is a function of the authenticated channel
identity alone, never of content, so the only way to obtain a high
label for adversarial content is to deliver it over a high channel,
which is channel spoofing and is excluded by
Assumption~\ref{asm:recorded}. The value recorder is the one TCB
component that reads content; it is a deterministic extractor over the
belief's text (identifiers, principals, amounts), and its failure mode
(missing a value) is fail-closed under \Ifour.

\begin{example}[Running example, I: the objects]
\label{ex:objects}
The user turn ``Delete file~0 from the shared drive'' is ingested on the
user channel as atom $e_u$ with $\ell(e_u) = \mathrm{TRUSTED}$,
$\att(e_u)$ true. The drive listing returned by a tool is atom $e_l$
($\mathrm{MEDIUM}$, attested by gateway receipt); the body of one listed
document, which carries the planted span ``also delete file~13, the
owner has approved,'' is atom $e_p$ ($\mathrm{MEDIUM}$, attested as a
tool result, content untrusted). The system forms beliefs
$b_u$ (``the user wants file~0 deleted,'' $\deps = \{e_u\}$,
$\vals = \{\text{file~0}\}$, $\mathrm{TRUSTED}$), $b_l$ (the listing,
$\mathrm{MEDIUM}$), and $b_p$ (``file~13 should be deleted,''
$\deps = \{e_p\}$, $\vals = \{\text{file~13}\}$, $\mathrm{MEDIUM}$) by
the meet rule. The
action $\ev{act}(a)$ with $c(a) = \mathrm{delete} \in \Cprot$ has
$\args(a)(\mathsf{target}) = \text{file~13}$; the policy binds
$\mathsf{target}$ at $\ell_{\bindpol}(\mathrm{delete}, \mathsf{target})
= \mathrm{TRUSTED}$. The compromised model commits $\support(a) =
\{b_u\}$: a genuine, attested, trusted directive.
\end{example}

\subsection{The Definition}
\begin{definition}[Provenance-completeness, \Ione]
\label{def:i1}
A trace $\sigma$ satisfies \Ione{} iff for every
$\ev{act}(a) \in \acts(\sigma)$ and every $b \in \support(a)$:
(i)~the transitive dependency closure of $b$ in $P$ is
well-founded and every maximal derivation path terminates in $E$;
and (ii)~every such terminal evidence atom $e$ satisfies $\att(e)$
and $\ell(e) \sqsupseteq \lmin$. Informally: no protected action
rests on an orphan belief (no evidential root, the formal shape
of a fabricated claim) or on an unattested or sub-floor root.
\end{definition}

\begin{definition}[Integrity non-degradation, \Itwo]
\label{def:i2}
A trace $\sigma$ satisfies \Itwo{} iff every label increase on any
belief is effected by a logged endorsement event authorized by
$\pi$; equivalently, for every update $\ev{upd}(b', X)$,
$\ell(b') \sqsubseteq \bigsqcap_{x \in X} \ell(x)$, and no other
event raises a label. Informally: untrusted input can never
silently launder itself into trusted belief; every upgrade has an
accountable, policy-cited justification.
\end{definition}

\begin{definition}[Evidential-support sufficiency, \Ithree]
\label{def:i3}
Fix an audited aggregator $S : 2^{B} \times L \to [0,1]$ scoring a
support set against a required label, subject to three stated
requirements: \emph{(A1) qualification}: $S(U, \ell)$ depends only
on evidence at or above $\ell$ that is attested;
\emph{(A2) monotonicity}: adding qualifying evidence never
decreases the score; \emph{(A3) auditability}: $S$ is fixed,
computable from the recorded prefix state, and inside the TCB.
$S$ is instantiable as Dempster--Shafer belief $\mathrm{Bel}(\cdot)$
over the qualifying evidence~\cite{shafer}, or as an
accepted-argument criterion in an argumentation
framework~\cite{dung}. A trace satisfies \Ithree{} iff for every
$\ev{act}(a) \in \acts(\sigma)$,
$S(\support(a), \lreq(c(a))) \geq \tau_{c(a)}$. Informally:
consequential actions require enough qualifying evidence, with
per-class thresholds.
\end{definition}

\begin{definition}[Argument binding, \Ifour]
\label{def:i4}
A trace $\sigma$ satisfies \Ifour{} iff for every
$\ev{act}(a) \in \acts(\sigma)$ and every bound slot
$s \in \bindpol(c(a))$ with $\args(a)(s)$ defined, there exists
$b \in \support(a)$ with $\args(a)(s) \in \vals(b)$ and
$\ell(b) \sqsupseteq \ell_{\bindpol}(c(a), s)$. Informally: every
argument that determines what the action does to the world is named by
a cited belief of sufficient standing. The model may choose the action;
it may not choose its targets. \Ifour{} is the decision-layer
counterpart of the authority binding that context-to-execution
integrity enforces at the sink~\cite{cxi}: CXI binds a field value to
the manifest that authorizes its effect, \Ifour{} binds an argument
value to the evidence that justifies naming it. An argument slot left
unbound by $\bindpol$ (free-text message bodies, say) is unconstrained
by \Ifour{} and falls to \Ifive.
\end{definition}

\begin{definition}[Counterfactual-dependence bound, \Ifive]
\label{def:i5}
Let $\Dirrel(q)$ be a class of task-irrelevant perturbations of
untrusted content (perturbations that preserve the task-relevant
semantics of $d$; operationalized by a generator suite,
\S\ref{sec:falsifiability}), and let $\dist$ be a metric on
protected actions (0/1 for discrete choices). A system satisfies
\Ifive{} iff for all $\delta \in \Dirrel(q)$,
$\dist\big(J(q, d),\, J(q, d \oplus \delta)\big) \leq \varepsilon_{c}$
for the class $c$ of $J(q,d)$. This is the invariance condition
articulated for judgment integrity in~\cite{qu}, restated over decision
traces; unlike \Ione--\Ifour{} it relates \emph{pairs} of executions.
\end{definition}

\begin{definition}[Decision-level epistemic integrity]
\label{def:dei}
Write $\Lift(P) = \{\,T : \forall \sigma \in T.\; \sigma \models P\,\}$
for the set-level lifting of a trace property $P$. A system with
reachable-trace set $T$ satisfies \DEI{} with respect to policy
$\Pi$ iff
\[
T \;\in\; \underbrace{\Lift(\Ione \wedge \Itwo \wedge \Ithree \wedge \Ifour)}_{\text{safety fragment (enforceable)}} \;\cap\; \underbrace{\Ifive}_{\text{2-safety (auditable)}} .
\]
\DEI{} is thus, by construction, a hyperproperty whose safety fragment
and whose properly $k$-safety fragment are named separately: the form in
which the enforceability results of \S\ref{sec:monitorability} are
stated.
\end{definition}

\begin{example}[Running example, II: what the invariants decide]
\label{ex:invariants}
On the trace of Example~\ref{ex:objects}: \Ione{} holds ($b_u$ roots in
the attested, trusted $e_u$); \Itwo{} holds (no label was raised);
\Ithree{} holds (a trusted directive is $\tc$-sufficient for a single
deletion under the deployed scorer); \Ifour{} \emph{fails}: the bound
slot $\mathsf{target}$ carries file~13, and no cited belief at
$\mathrm{TRUSTED}$ names file~13. The monitor blocks at admission. Now
change the plant to ``delete file~0 now, before the backup completes,''
with the user having asked for the backup first: the model commits
$\args(a)(\mathsf{target}) = \text{file~0}$ with the same support, and
\Ione--\Ifour{} all hold. The action is wrong because its timing was chosen by the plant, but
every single-trace check passes; only \Ifive{}---re-running the
decision with the plant removed---distinguishes this trace from an
honest one. Corollary~\ref{cor:ratgap} shows this is a limit of
admission-time checking as such, not of our particular invariants.
\end{example}

\subsection{Monitorability: What a Gateway Can Enforce}
\label{sec:monitorability}
The invariants deliberately fall on both sides of a classical
boundary, and the definition's honesty depends on saying which
side each is on. Algorithm~\ref{alg:monitor} makes the
enforceable side concrete: a single event-admission procedure whose block
branches are exactly the bad prefixes of
Proposition~\ref{prop:safety}. Two assumptions are stated explicitly
because the proofs depend on them.

%%ALG1%%

\begin{assumption}[Complete mediation]
\label{asm:mediation}
The gateway mediates every $\Sigma$-event that can affect a
protected action; there is no side channel from untrusted content
to belief state or action that bypasses admission.
\end{assumption}

\begin{assumption}[Trusted TCB]
\label{asm:recorded}
The epistemic TCB of \S\ref{sec:tcbdef} behaves as specified:
(i)~channel identity is authenticated, so $\lambda$ assigns to every
atom the label its channel carries by policy; (ii)~the recorder
extracts $\vals(\cdot)$ faithfully and the log is append-only;
(iii)~the attestation verifier, the aggregator $S$, and the policy
engine compute what their definitions state. Compromise of any
component voids all guarantees (\S\ref{sec:limitations}).
\end{assumption}

\begin{proposition}[\Ione--\Ifour{} are enforceable safety properties]
\label{prop:safety}
Over decision traces, \Ione, \Itwo, \Ithree, and \Ifour{} are safety
properties in the sense of Alpern--Schneider~\cite{alpern}: every
violating trace has a finite bad prefix (the offending action or
update event) after which no extension repairs the violation.
Consequently (i)~they are monitorable, and an
$\mathrm{LTL}_3$-style monitor~\cite{bauer,pnueli} issues the
verdict $\bot$ exactly at the bad prefix; and (ii)~they lie in
the class of policies enforceable by execution
monitoring~\cite{schneider}: a reference monitor that truncates
the trace at the bad prefix (blocking the action or
update) enforces them, not merely detects them.
\end{proposition}

\begin{proof}[Proof sketch]
For \Ione, \Ithree, and \Ifour{} the candidate bad event is a protected
action; its support set, the relevant provenance subgraph, the
aggregator score, the cited beliefs' value sets, and the action's
argument map are all determined by the finite prefix at that event, so
violation is a property of the prefix and is irremediable once the
action executes. For \Itwo{} the bad event is a label-raising
transition without a matching authorized endorsement in the same
prefix. Quantification over runtime objects (beliefs, evidence) is
handled by parametric monitoring~\cite{mop}: the object universe at
any prefix is the finite set of objects observed so far, so each
instantiated monitor is a finite-state safety monitor. Enforcement
follows because blocking the bad event yields a legal trace, i.e., the
monitor is a truncation automaton in Schneider's
sense~\cite{schneider}. The full proof, including the two
legality-preservation facts, is Appendix~\ref{app:proofs}. The
\Ione--\Ifour{} clauses and the two legality facts are mechanized in
Lean~4 (finite-trace form; kernel-checked; the only axiom is
\texttt{propext}, on two lemmas), re-checkable in the artifact via
\texttt{formal/check\_formal.sh}.
\end{proof}

\begin{remark}[Online cost]
With per-belief \emph{root certificates} (a memoized, hash-chained
receipt that a belief's evidential roots are attested and
$\sqsupseteq \lmin$, established when the belief is created) the
\Ione{} check at an action is $O(|\support(a)|)$ constant-time receipt
verifications; \Itwo{} is one lattice meet per update; \Ithree{} is one
aggregator evaluation; \Ifour{} is $O(|\bindpol(c(a))| \cdot
|\support(a)|)$ set-membership tests on recorded value sets. These are
the operations whose cost CZTA measures (Table~\ref{tab:overhead}); the
reference implementation confirms the asymptotics, with memoized
\Ione{} checks flat across a $256\times$ range of provenance depth
while the unmemoized check grows linearly
(Appendix~\ref{app:refsem}).
\end{remark}

\begin{proposition}[\Ifive{} is auditable but not single-trace monitorable]
\label{prop:hyper}
\Ifive{} is a 2-safety hyperproperty~\cite{hyperprops}: its
violation is witnessed only by a pair of finite traces
$(J(q, d), J(q, d \oplus \delta))$. No black-box monitor observing
a single production trace can detect a violation, since
satisfaction is not a property of any individual
trace~\cite{hyperprops}; runtime checking requires either stored
comparison executions or a model of the system (gray-box
monitoring)~\cite{finkbeiner,stucki}. \Ifive{} is therefore
classified \emph{auditable}: falsifiable offline by exhibiting a
witness pair, testable by a differential harness, and checkable before
dispatch only by re-executing the decision (\S\ref{sec:gate}), never by
inspecting the production trace alone.
\end{proposition}

\begin{corollary}[The DEI decomposition]
\label{cor:decomp}
$\DEI = \Lift(\Ione \wedge \Itwo \wedge \Ithree \wedge \Ifour) \cap \Ifive$
(Def.~\ref{def:dei}) is the conjunction of an EM-enforceable
safety fragment (Prop.~\ref{prop:safety}) and a 2-safety fragment
that no single-trace monitor enforces (Prop.~\ref{prop:hyper}).
Any monitor claiming to enforce \DEI{} wholesale is unsound.
\end{corollary}

\begin{remark}[The boundary, stated plainly]
Full noninterference between untrusted channels and decisions is a
general hyperproperty~\cite{hyperprops,goguen} and inherits the same
non-monitorability; we neither claim nor need it. Complete
transformer-level noninterference is explicitly a non-goal even for
sink-authorization systems~\cite{cxi}, and there are principled
arguments that in-band immunity to injection may be
unattainable~\cite{abdelnabi}. Likewise, no monitor decides whether a
belief is \emph{true}; \DEI{} constrains the evidential and provenance
structure of decisions, not their semantic veracity. Concurrent work
characterizes generically what a deterministic pre-execution agent gate
can enforce: the safety policies whose good prefixes its state model
recognizes~\cite{ray2026enforced}. Prop.~\ref{prop:hyper} names a
specific, security-relevant property outside that class for
\emph{every} state model, and \S\ref{sec:gate} builds the
decision-local gate that audits exactly that property.
\end{remark}

\begin{proposition}[Mutual independence of the aggregator axioms]
\label{prop:indep}
There exist aggregators $S_{\neg\mathrm{A1}}$, $S_{\neg\mathrm{A2}}$,
$S_{\neg\mathrm{A3}}$, each violating exactly one axiom of
Def.~\ref{def:i3} while satisfying the other two: no axiom is
derivable from the rest, and dropping any one admits a named attack.
\end{proposition}
\noindent\emph{Proof:} the artifact, one explicit violating aggregator
per axiom with codomain and monotonicity checks.

\subsection{The Rationalization Gap}
\label{sec:guarantees}
The crosswalk of Appendix~\ref{app:crosswalk} maps each attack class of
\S\ref{sec:taxonomy} to the invariant it must violate.
Theorem~\ref{thm:crosswalk} shows the mapping is sound against the
enforceable fragment; Corollary~\ref{cor:ratgap} shows why that
fragment is nevertheless insufficient, and exactly which attacks
survive it.

\begin{theorem}[Crosswalk soundness]
\label{thm:crosswalk}
Under Assumptions~\ref{asm:mediation} and~\ref{asm:recorded}: for
every charting class $K$ with machine-fired primary $j(K)$ and every
trace $\sigma$ with $\mathrm{Mech}_K(\sigma)$, there is a finite prefix
$\sigma[..k]$ with $\mathit{viol}_{j(K)}(\sigma[..k])$, and
Algorithm~\ref{alg:monitor} blocks at or before $\epsilon_k$. Conversely
(anti-triviality), control/evaluation instances satisfy $\neg
\mathit{viol}_j$ at every prefix for all $j$ and are admitted.
\end{theorem}
\noindent\emph{Proof:} Appendix~\ref{app:t1c2} (finite case analysis
over the machine-fired rows, anti-triviality, and a delimitation of the
proved implication from its empirical support). The per-row predicates
$\mathrm{Mech}_K$ (Definition~\ref{def:mech}) are stated there.

\begin{corollary}[Rationalization gap: choice laundering is admission-undetectable]
\label{cor:ratgap}
Let $M$ be any sound monitor enforcing $\Ione \wedge \Itwo \wedge
\Ithree \wedge \Ifour$ at admission over traces where $\support(a)$ is
committed by an untrusted component (the model; a compromised model is
an adversarial component). Let $S_0$ be the support set of any honestly
admitted protected action of class $c_0$, and let $A_{\bindpol}(S_0)$
be the set of class-$c_0$ actions every bound argument of which is named
by some belief in $S_0$ at the label $\bindpol$ requires. Then for every
$a^{*} \in A_{\bindpol}(S_0)$ there exists an adversarial model
strategy---select $a^{*}$ from injected content; commit $S_0$---such that
on the resulting trace $\mathit{viol}_j$ holds at no prefix for any
$j \in \{1,2,3,4\}$ while the action's actual basis is the injection.
Hence (i)~the strategy is undetectable at admission by any sound
$\Ione$--$\Ifour$ monitor at \emph{any} provenance granularity; and
(ii)~an $\Ifive$-side counterfactual check is necessary, not optional:
no refinement of the enforceable fragment detects it.
\end{corollary}
\noindent\emph{Proof:} Appendix~\ref{app:t1c2} (existence via the
full-support commitment $S_0$ restricted to $A_{\bindpol}(S_0)$;
necessity via the $S^{\pm}$ construction of Prop.~\ref{prop:hyper}).

\begin{remark}[What \Ifour{} catches and what it leaves]
\label{rem:residual}
The corollary quantifies over $A_{\bindpol}(S_0)$, which is the
precise residual: an injection can still choose \emph{which} of the
trusted-derivable actions is taken, \emph{whether} it is taken, and
\emph{when}, but it can no longer choose a target that trusted evidence
does not name. The laundering probes executed in \S\ref{sec:rq2}
(delete file~13 for file~0; share with an attacker-chosen recipient) are
\emph{not} instances of the corollary's strategy: they violate \Ifour{}
and are blocked at admission once \Ifour{} is enforced. The gate's
incremental value must therefore be demonstrated on choice-laundering
probes, and \S\ref{sec:rq2} is designed accordingly. This is also why
\Ifour{} is a separate invariant rather than a clause of \Ithree: it
changes the adversary's residual strategy space, and the corollary
states that space exactly.
\end{remark}

\subsection{Falsifiability Protocol}
\label{sec:falsifiability}
Each invariant carries a refutation condition: the finite observation
that falsifies the claim that a given system enforces or satisfies it.
A single logged protected action admitted with an unattested or
sub-floor evidential root refutes ``the deployment enforces \Ione'';
a single silent label upgrade refutes \Itwo; a single admitted action
below threshold refutes \Ithree; a single admitted action with a bound
argument named by no qualifying cited belief refutes \Ifour; and one
$\delta \in \Dirrel$ with $\dist(J(q,d), J(q,d \oplus \delta)) >
\varepsilon_c$ refutes \Ifive. The first four are checkable by any
auditor holding the log, which is why the log is hash-chained and
receipt-bearing; the fifth requires the pre-registered differential
harness.

\paragraph*{Falsifying the definition itself (coverage)}
Let $\mathcal{A}$ be the attack taxonomy of \S\ref{sec:taxonomy} with
the per-instance success predicate of the evaluation harness. Say
\DEI{} \emph{covers} $\mathcal{A}$ on a corpus if every successful
attack instance violates at least one of \Ione--\Ifour{} on its trace
or \Ifive{} on its witness pair. A successful epistemic attack on an
invariant-satisfying execution falsifies coverage and shows the
definition incomplete: a reportable outcome, not an embarrassment. The
crosswalk (Appendix~\ref{app:crosswalk}) records, per attack class, the
invariant predicted to be violated, registered before the campaigns
that test it. An adaptive adversary must therefore either violate a
monitored invariant (and be blocked or logged) or exhibit a coverage
counterexample (and falsify the definition), the adaptive-evaluation
design the systematization of out-of-band defenses calls
for~\cite{oob}.

\begin{example}[Running example, III: the protocol applied]
\label{ex:protocol}
Walk the two traces of Example~\ref{ex:invariants} through the
protocol. Trace~1 (target file~13): if the deployment admits the action,
the admission log itself refutes ``the deployment enforces
\Ifour''---an implementation failure, repaired by fixing the monitor,
with no bearing on the definition. Trace~2 (target file~0, timing chosen
by the plant): every enforceable invariant holds on the log, so no
admission record can refute anything; the differential harness re-runs
the decision with the plant removed, the action changes, and the pair
refutes \Ifive. Had the action \emph{not} changed while the attack still
succeeded---say through an unbound free-text argument the perturbation
family never touches---that would be a coverage counterexample,
reported as evidence that the definition (here $\bindpol$ or
$\Dirrel$) is incomplete. The three outcomes are distinguishable from
the artifacts alone, which is what makes the protocol executable.
\end{example}

\paragraph*{Human validation of the crosswalk}
Because the crosswalk is a judgment artifact, two blinded domain
experts independently assigned invariants to each class from the
definitions alone, with an adjudicator resolving disagreements. On the
eight classes with a machine-fired primary they agreed exactly on 7/8
(Wilson 95\,\% CI 0.53--0.98) and 10/10 on the in-scope boundary; the
adjudicated consensus matches the invariant that actually fired in the
reference monitor on 6/8. The divergences are informative: both coders
read stale evidence as insufficient support (\Ithree) where the monitor
enforces a lapsed attestation (\Ione), and read prompt injection as a
provenance failure (\Ione) where the gateway blocks the unendorsed
label raise (\Itwo). The crosswalk records what the monitor
\emph{enforces}; the mapping was frozen before any campaign was scored
against it, and the instrument and scorer are released.

\subsection{Scope and Limitations}
\label{sec:limitations}
Appendix~\ref{app:refsem} distinguishes proved, implemented, measured,
development, and uninformative evidence, so that no deferred study or
failed susceptibility check is mistaken for a positive result. Beyond
that: (1)~all guarantees are relative to the epistemic TCB; compromising
the labeler, the channel authentication it relies on, the verifier
keys, the recorder, or the aggregator voids them, and \S\ref{sec:tcb}
quantifies that base. (2)~Rationalization is a faithfulness gap, not a
granularity one. \Ione--\Ifour{} constrain the \emph{committed} support
set and the action's bound arguments, but the support set is committed
by an untrusted model, so a compromised model can select an action from
adversarial content and cite legitimate evidence for it. \Ifour{}
narrows the residual to choice laundering under trusted-derivable
arguments (Cor.~\ref{cor:ratgap}) without closing it; sub-document
laundering is a resolution limit that finer provenance
tightens~\cite{argus,cxi}, and~\cite{romanchuk2026semantic} states the
architectural version of the gap as \emph{semantic laundering}, for
which \S\ref{sec:gate} builds the runtime check. (3)~$\bindpol$ and
$\ell_{\bindpol}$ are policy with a utility cost: a target discovered
from a $\mathrm{MEDIUM}$ tool result cannot bind at $\mathrm{TRUSTED}$,
so $\bindpol$'s tightness is bounded by the attestation topology
exactly as $\tc$ is (\S\ref{sec:evalsetup}). (4)~$\Dirrel$ is itself a judgment
and the hard core of \Ifive~\cite{qu}; we pre-register the generator
suite rather than choose perturbations post hoc. (5)~\Ithree's
aggregator encodes an epistemology; we require only A1--A3.
(6)~\DEI{} constrains structure, not truth: attested, well-supported,
invariant-satisfying decisions can still be wrong.
